% Page 035

\seite{35}

\margmark{Übergangsstelle\\dem alten und neuen\\Pfarrstellen.}%
Abschied ist allen begonnen wir, es\ob
müßte am großen bei Eingenommen\ob
um Gunst wirklich eine Bismarckrede\unsicher{} erregen\ob
worden.

\pend
\begin{verse}
Seffern, den 11.\ März\unsicher\\
Ettges. {[Unterschrift]}
\end{verse}
\pstart

\pend
\abschnitt{Osterprüfung}
\pstart

Am 11.\ März fand durch Herrn mit\ob
Ortsschulinspector Ettges die übliche Oster\ohb
prüfung statt.

\pend
\jahresueberschrift{Das Schuljahr 1903}
\pstart

Es wurden entlassen zu Ostern 1903,\ob
3 Knaben und 3 Mädchen. Die Knaben\ob
widmen sich der Landwirthschaft. Aufgenommen\ob
wurden fünf Knaben. Alle Kinder sind\ob
katholisch und willfährig.

\pend
\abschnitt{Vertretung}
\pstart

Auch in Einberufung des Herrn Lehrer Fritz\ob
nach Seffern, zu einer militärischen\ob
Übung in Frankfurt, wurde der hiesige\ob
Lehrer mit der Vertretung betraut. Und es\ob
fand abwechselnd Unterricht in beiden Schu\ohb
len statt.

\pend
\abschnitt{Schulfeien}
\pstart

Am 2.\ September fand die Feier des Sedans\ohb
festes statt. Einverständnis mit dem Vikar und\ob
Familien nach einige Lieder warten anbei\ob
östlichen Bräuche Schwimmen der, und\ob
Abschließend wurde im Stillen, an und mit\ob
in der Amts und beim Kränze, und im Stillen\ob
als Frankreich gefeiert. Ein stets Pflicht und\ob
einem Gruß auf unsern Kaiser.
\pend
